\documentclass[11pt,oneside]{book}

\usepackage[T1]{fontenc}
\usepackage[utf8]{inputenc}
\usepackage{lmodern}
\usepackage[a5paper,margin=25mm]{geometry}
\usepackage[poemsintoc,numberpoems,clearpageafterpoem]{poetrytex}

\renewcommand*{\pttitle}{Small Weather}
\renewcommand*{\ptsubtitle}{An original field anthology}
\renewcommand*{\ptauthor}{A. Rowan}
\renewcommand*{\ptdate}{2026}
\renewcommand*{\ptdedication}{For the people who notice the sky}

\begin{document}

\maketitle
\makededication[center]
\maketoc
\maketop

\poemgroup{First Light}

\begin{poem}{Window Before Breakfast}{A. Rowan}
The kettle gathers a small cloud,\\
and the glass takes it for weather.\\
Outside, one bicycle turns the corner\\
with a bright red lamp still burning.\\
I write the day down softly,\\
so the morning has room to arrive.
\end{poem}

\begin{poem}{Map of the Quiet Street}{A. Rowan}
At the crossing, the signal changes\\
before anyone asks it to.\\
A plane draws a white measure\\
between two roofs.\\
The longest route home\\
is the one that keeps its promises.
\end{poem}

\begin{annotation}{Notes on First Light}
The first group uses short, line-oriented poems and a shared author line. In a longer anthology, replace the group title, poem title, and subtitle arguments while retaining the same page and list-of-poems structure.
\end{annotation}

\poemgroup{After Rain}

\begin{poem}{Puddles Remember}{A. Rowan}
Rain leaves a second city\\
under the first one: \\
steps become bridges,\\
curbs become shorelines,\\
and every passing cloud\\
is carried home in a shoe.
\end{poem}

\begin{poem}[center]{The Last Drop}{A. Rowan}
One drop waits on the railing.\\
It holds the garden upside down.\\
When it falls, the garden returns\\
to being only a garden.
\end{poem}

\begin{annotation}{Notes on After Rain}
Annotations are separate prose sections connected to a group of poems. They can hold editorial notes, translations, source notes, or a short reflection without changing the poem layout itself.
\end{annotation}

\end{document}
